\chapter{Overall summary and conclusions}\label{sec:6}

This book had two major aims. The first aim was to probe a compositional analysis of meaning change of modal indefinites. In order to achieve this goal, it was necessary to characterise and define different meanings of \textit{cualquiera} in Modern Spanish varieties. The reason why \textit{cualquiera} is particularly interesting for a linguistic study is that this word has different meanings that appear to be unrelated such as the Free Choice (FC) meaning ‘any’ and evaluative meanings such as ‘ordinary’ and ‘low value’ (which I called EVAL1) and ‘nonsense’ (EVAL2). I focused on European and Argentinian Spanish, because these varieties show different degrees of lexicalisation of these meanings. One research question was whether the different meanings of \textit{cualquiera} in Spanish varieties represent a case of polysemy (Hypothesis 1) or whether \textit{cualquiera} has only one basic meaning (assumedly the FC meaning ‘any’) and other meanings can be inferred from the general meaning pragmatically (Hypothesis 2). Moreover, I investigated the relation between syntax and semantics, i.e. whether certain meanings appear under particular morphosyntactic conditions, which is the core assumption of the compositional framework of meaning change. The second aim of this book was to investigate the diachrony of these meanings, i.e. when 
they appeared over time and under which linguistic conditions they occurred. The major scope of this book is thus a synchronic and diachronic description and analysis of different meanings of \textit{cualquiera} in European and Argentinian Spanish. In the first half of this book, I investigated the synchrony of \textit{cualquiera}, and in the second half I investigated its diachrony. I will summarise how I proceeded in my analysis and what the most important results and new findings to the study of \textit{cualquiera} are.

  In order to describe the synchronic meanings of \textit{cualquiera}, I first reviewed the data and various analyses as discussed in the literature in Modern European and Argentinian Spanish. I have shown that semantic and syntactic analyses are missing on \textit{cualquiera} with EVAL1 ‘unremarkable’ present in European and Argentinian Spanish and EVAL2 ‘nonsense’ present only in Argentinian Spanish.

  In \chapref{sec:3}, I presented new data and a new analysis of N \textit{cualquiera} with EVAL1. I analysed new data from a subcorpus of Corpus del Español (CdE\_web) and have identified various flavours of the evaluative meaning ‘unremarkable’, namely COM (common: ‘regular, normal’), DER (derogative: ‘ordinary, without any particular quality worth mentioning’). I argued that EVAL1 in European Spanish is a special case of interpreting \textit{cualquiera} as a quantifier with a free choice interpretation ‘any’. I have shown that by deriving the evaluative meaning from the quantifier meaning, \textit{cualquiera} is not a case of polysemy (i.e. Hypothesis 2 of \sectref{sec:1.1} is verified and Hypothesis 2 is falsified, contra \citealt{AlonsoMenendez2018}). In order to derive the evaluative meaning ‘ordinary’ from ‘any’, I have analysed \textit{cualquiera} as a quantificational predicate (modifier) similar to adjectives, which has the interpretation of a comparative clause expressing similarity between the actual referent and a comparison class, i.e. \textit{Juan es un hombre cualquiera} is interpreted as ‘Juan is a man and he is like anyone else’ (whereby ‘anyone else’ is a set of comparable male individuals, see \sectref{sec:3.3}). The similarity between individuals entails the meaning ‘common’ as having similar properties like anyone else entails having common properties. The derogatory meaning of ‘low value’ or ‘nothing special’ is a pragmatic inference of the ‘common’ meaning. The free choice meaning ‘any’ is the basic lexical meaning of \textit{cualquiera} and EVAL1 ‘ordinary’ follows from this basic meaning. This analysis is the first analysis that provides an explanation how different meanings and functions of \textit{cualquiera} are interconnected and it provides further support for the compositionality view on meaning change.

  I discussed a small data set in the corpus CdE\_web (e.g., definite determiner N \textit{cualquiera}) that shows signs of lexicalisation of the evaluative meaning. According to this data set, the SIM meaning (i.e. \textit{cualquiera} as ‘like anyone else’) has been reanalysed as a property of kinds of individuals with the meaning ‘common/typical/ordinary’. This lexical meaning co-occurs with definite nouns and quantified nouns (e.g. \textit{el hombre cualquiera} ‘the typical person/man’ or \textit{algún día cualquiera} ‘on some ordinary day’) as well as with coordination (e.g. \textit{un ser mediocre y cualquiera} ‘a mediocre and ordinary human being’).

  In \chapref{sec:4}, I analysed \textit{cualquiera} with EVAL1 and EVAL2 in Argentinian Spanish and suggested an analysis for the interpretations and licensing of \textit{cualquiera} in episodic contexts in this variety. Episodic contexts usually do not license modal indefinites like \textit{cualquier} ‘any’. I suggested that \textit{cualquiera} with EVAL1 should be analysed as a degree adjective in Argentinian Spanish. I have argued that the difference between European Spanish and Argentinian Spanish \textit{cualquiera} with EVAL1 is primarily syntactic: In Argentinian Spanish but not in European Spanish, \textit{cualquiera} is used as a predicative adjective with degree modification (e.g. \textit{es muy cualquiera} ‘he/she/it is very ordinary’). I analysed the meaning of the adjective \textit{cualquiera} as ‘bad’, as a pejorative interpretation of \textit{cualquiera} as ‘common’ or ‘ordinary’ under the assumption that speakers who use the pejorative version ‘bad’ assume that ‘being ordinary’ counts as a negative property in contrast to ‘being special’ or ‘having a distinctive property’. I have proposed analysing EVAL2 ‘nonsense’ in parallel with the free choice indefinite \textit{cualquier cosa} ‘anything’, which is interpreted in episodic contexts as ‘anything possible’ or as ‘something random’. The evaluative meaning ‘nonsense’ is a pragmatic inference, which comes about as follows. The judge (usually the speaker) evaluates the modal component of the free choice \textit{cualquiera} (‘every possibility is an option’) on a scale of preferences, namely as the least preferred option. Under verbs of saying like \textit{decir} or \textit{mandar} ‘to say’ in Argentinian Spanish, this pragmatic inference works as follows. The expression in Argentinian Spanish \textit{decir/mandar cualquiera} ‘saying anything possible’ implies saying contradictory things or saying something random or saying unreflected things according to the judge (usually the speaker). This pragmatic evaluation has been lexicalised in Argentinian Spanish under transitive verbs of saying and this evaluative meaning further generalised to other verbs, e.g. predicative verbs (e.g. ArgSp \textit{lo que me decís es cualquiera} ‘what you are telling me is bullshit/nonsensical’). The diachronic data of Argentinian Spanish \textit{cualquiera} is unfortunately too scarce to draw strong conclusions. Nevertheless, there are some hints that suggest that the origin of the meaning ‘nonsense’ is its occurrence under verbs of \textit{decir} ‘to say’. Concerning the research question of this book whether EVAL1 and EVAL2 are lexicalised meanings, I have argued that the answer is positive for Argentinian Spanish (or, in other words, Hypothesis 1 holds) as these meanings correlate with different syntactic categories than the indefinite \textit{cualquiera} with FC ‘anything’, namely degree adjectives and evaluative nouns. However, what looks like a case of polysemy in Modern Argentinian Spanish is not a polysemy from the diachronic perspective. I argued that, in diachrony EVAL1 and EVAL2, are the result of pragmatic inferencing that has been lexicalised and has driven syntactic reanalysis into different categories in Modern Argentinian Spanish, namely from a pronoun into gradable adjective with the meanings ‘ordinary’ and ‘nonsensical’ and into an evaluative noun with the meaning ‘nonsense’. I have shown that this path of reanalysis from FC \textit{cualquiera} into evaluative \textit{cualquiera} is not a special property of \textit{cualquiera} in Argentinian Spanish, but is also present in French (e.g. \textit{n’importe quoi}), Mexican Spanish \textit{equis} (‘any’ in modal contexts and ’ordinary’ under predicative verbs) \citep{Kellert2020semantic} and other, non-Romance languages, as for instance in Czech (see \citealt{Simik2008}). \chapref{sec:4} presents the first in-depth contrastive analysis of \textit{cualquiera} in Argentinian and European Spanish, based on corpus data and quantitative analysis.

  After the synchronic analysis of evaluative meanings and the diachronic analysis of EVAL2 in Argentinian Spanish, I focused on the diachronic analysis of \textit{cualquiera} in European Spanish in order to find evidence for my synchronic analysis, namely that the free choice meaning is the basic meaning and the evaluative meanings are secondary or following from the free choice meaning. My assumption was that the free choice meaning occurred diachronically first and the evaluative meanings later. Moreover, I looked at the change of linguistic conditions (e.g. change of episodic contexts, syntactic position, etc.) to see how the meanings arise diachronically.

  In \chapref{sec:5}, I have shown on the basis of quantitative data and their statistical analysis on the basis of historical data from CdE when and how \textit{cualquiera} has developed its functions and meanings in the history of Old and Middle Spanish as well as Modern European Spanish. I analysed different syntactic properties such as the position of \textit{cualquiera} with respect to the verb and the noun and have shown on the basis of statistical tests that \textit{cualquiera} has been used increasingly in different syntactic structures such as declaratives, in postverbal position, etc., which I interpreted as an argument for the grammaticalisation hypothesis. I have modelled the path of grammaticalisation of \textit{cualquiera} in a generative framework and have analysed the meaning components that have changed from the relative clause with free choice interpretation \textit{qual quier,e,a} ‘which one wants’ to \textit{qualquier,-e,-a} ‘whichever’. The most important semantic component that changed from the meaning of relative clause \textit{qual quier,e,a} to the free choice indefinite \textit{qualquier,-e,-a} is related to the judge parameter. The free choice interpretation of the indefinite (in contrast to the relative clause) is no longer dependent on the generic \textit{pro} as in \textit{qual pro\textsubscript{generic}} \textit{quier,-e,-a,} but can be associated with the agent of the matrix clause, (e.g. \textit{Pablo} \textit{puede invitar a cualquiera} ‘Pablo can invite anyone’). What both linguistic expressions (relative clause Old Spanish \textit{qual quier,-e,-a} and the modern indefinite \textit{cualquiera}) have in common is that they do not contribute to the truth conditional semantic meaning of the clause ‘Pablo can invite someone’. They have rather the function of a side-comment, e.g. ‘Pablo has free choice of inviting whoever he likes’. I have shown diachronic evidence for the side-comment analysis as relative clauses such as \textit{qual quier,-e,-a} behave very much like adjuncts and they can even interrupt a sentence (e.g. \textit{pon una estrella fixa qual quisieres por su longura del eguador} ‘put a star [which you want] along the equator’).

  \largerpage
  Another important finding of my diachronic study is that the postnominal \textit{cualquiera} is the result of a process by which the relative clause \textit{qual} \textit{quier,e,a} has been lexicalised into a single compound, {meaning that the lexicon has been updated with a new entry \textit{cualquiera} as a nominal modifier. In this function as a nominal modifier, it has persevered in its semantic function as a nominal property (Noun + \textit{qualquier,-e,-a}). The evaluative meaning of postnominal \textit{cualquiera}, which appeared in the 17{th} century, is thus directly related to the property function of relative clauses in a postnominal position in contrast to the prenominal quantifier \textit{cualquier}. I have concluded that postnominal \textit{cualquiera} did not undergo grammaticalisation into a quantificational determiner as the prenominal \textit{cualquier} did. It is thus necessary to distinguish between different syntactic functions of \textit{cualquier(a)} (prenominal determiner, postnominal modifier, pronoun) to understand the different processes involved in the diachrony of these elements. This chapter provides the first in-depth diachronic analysis that combines statistical analysis with a theoretical explanation of meaning change of \textit{cualquiera} and offers a new perspective on the grammaticalisation of \textit{cualquiera}. A more general contribution to the research on \textit{cualquiera} is that I have shown the necessity of looking at different syntactic properties and categories of \textit{cualquiera} in diachronic and diatopic variation to understand the processes involved in the semantic and syntactic changes.}

  \largerpage
  I have shown that there are (at least) two diachronic paths of \textit{cualquiera} depending on its part of speech and syntactic position,  which can be summarised as follows. For the diachronic development of postnominal \textit{cualquiera} in European Spanish and Argentinian Spanish the following steps have been identified:

  \begin{enumerate}[label=\arabic*., leftmargin=2em]
    \item Use of postnominal \textit{cualquiera} often after \textit{otro ‘other’ + noun} with a complement clause after \textit{quier} expressed by \textit{que ‘that’} and a selected subjunctive (e.g. \textit{qual quier que sea}) with modal interpretation ‘whichever one wants that it may be’ in 13th century Old Spanish:
    \begin{quote}
        \gll \textit{e.g.} \textit{o} \textit{de} \textit{otra} \textit{razon} \textit{qualquier} \textit{que} \textit{sea}\\
        {} or of another reason \textsc{qualquier} that may.be\\
        \glt ‘or for some other reason whichever one wants that it may be’ (13th c., Alfonso X, \textit{Siete Partidas})
    \end{quote}
    
    Use of nouns between \textit{qual} and \textit{quier}, which show that \textit{qual quier} could be analysed compositionally in 13th century \citep{Pato2012}:
    \begin{quote}
        \gll \textit{e.g.} \textit{[...]} \textit{en} \textit{qual} \textit{manera} \textit{quier} \textit{que} \textit{sean}\\
         {} {} \textit{in} \textit{which} \textit{manner} \textit{wants} \textit{that} \textit{be.\textsc{sbjv}}\\
        \glt ‘in whatever manner’ (13th c., Alfonso X, \textit{Libro de las cruces})
    \end{quote}

    \item Recategorisation of postnominal \textit{cualquiera} as one word, lexicalisation of the modal meaning (‘which one wants’ $\rightarrow$ ‘any’) (see also \citealt{Elvira2007}). Empirical evidence:
    \begin{enumerate}[label=\alph*), leftmargin=1.5em]
        \item disappearance of \textit{qual NP quier}, 
        \item orthographic non-separated form \textit{qualquier} increases, 
        \item decrease of complement clauses after \textit{quier} and selected subjunctive: \textit{qual quier que ...}, 
        \item increasing use of postnominal \textit{cualquiera} with an indefinite N as in \textit{\textsc{un} N cualquiera}.
    \end{enumerate}

    \item Use of \textit{cualquiera} in predicative and postnominal position under negation (e.g. \textit{no es un día cualquiera, es un día especial} ‘it’s not an ordinary day, it’s a special day’) (in 16th c. in European Spanish until today in European Spanish and in Argentinian Spanish).
    
    \item Increase of copular predicates without negation in present and past tense (e.g. \textit{es un día cualquiera} ‘it’s an ordinary day’) (from 18th c. in European Spanish).
    
    \item Extension of the use of \textit{\textsc{un} N cualquiera} to definite articles + N \textit{cualquiera} and quantified nouns as in \textit{la/mucha gente cualquiera} ‘the/many ordinary people’ (in European Spanish in 20th c. and Argentinian Spanish).

\end{enumerate}

The diachronic development of postnominal \textit{cualquiera} continued in Argentinian Spanish, but not in European Spanish:

\begin{enumerate}[label=\arabic*., leftmargin=2em, start=6]
    \item Use of \textit{cualquiera} as a predicative complement without the noun (in Argentinian Spanish in 19th c. until today), e.g. \textit{Su ocupación era cualquiera.} ‘His position was nothing important.’
    
    \item Appearance of degree modification with the predicative complement (in Argentinian Spanish in 20th c.) as in \textit{re/muy/pero/tan/medio cualquiera} ‘very/really/a bit ordinary’.
    
\end{enumerate}

The diachronic development of pronominal \textit{cualquiera} has been identified in Argentinian Spanish, but not in European Spanish:

\begin{enumerate}[label=\arabic*., leftmargin=2em]
    \item Use of the pronoun \textit{cualquiera} and \textit{cualquier cosa} with a modal meaning (‘any (thing)’) from Old Spanish until today in European Spanish and Argentinian Spanish.
    
    \item Increasing use of the pronoun \textit{cualquiera} and \textit{cualquier cosa} under a particular lexical verb group, namely verbs of saying with present and past tense morphology as in \chapref{sec:4} (in 19th c. in Argentinian Spanish until today, not in European Spanish).
    
    \item Extension of the use of verbs of saying to other verbs and appearance of nominal modification of \textit{cualquiera} and \textit{cualquier cosa}, e.g. \textit{cualquiera total} ‘complete bull-shit’ (in Argentinian Spanish in 20th c., not in European Spanish).
\end{enumerate}


{I have argued that the new evaluative meaning ‘ordinary’ in Argentinian Spanish. is the result of reanalysing the old meaning ‘any’ with contextual information
(\citealt{Blank1997,Traugott1989,Eckardt2006}, among others). More precisely, the free choice meaning (the original meaning) is enriched by the contextual information triggered as the result of using \textit{cualquiera} in the postnominal position of a negated predicative complement that is associated with focus or emphasis (e.g. \textit{no es} [\textit{un día cualquiera}\textsubscript{focus}]\textsubscript{predicative complement}}). This focus or emphasis of the postnominal position of a predicative complement used under negation expresses a contrast between free choice alternatives of \textit{cualquiera} and another focus alternative, which is the opposite of free choice, namely an alternative that can be described as specific or definite and which is recoverable from context (e.g. \textit{no es un día cualquiera, es un día especial} ‘it's not just any day, it's a special day’). This contrast between free choice alternatives and the specific contextual alternative has in turn induced a contrast of properties between \textit{cualquiera} characterising free choice alternatives as ‘non-specific’, ‘general’ or ‘common’ as opposed to ‘specific’, ‘unique’ or ‘uncommon’, which becomes clearly visible in the use of definite articles and quantified nouns with [N \textit{cualquiera}] (e.g. \textit{la gente cualquiera} ‘the common people’). I adapt the new meaning as ‘common’ following the semantic literature on predicate modification of kind nouns (\citealt{Carlson1977,GehrkeMcNally2015}) to make the new meaning of \textit{cualquiera} more explicit in formal semantic terms. In my reconstruction of the reanalysis path, the change from a quantificational indefinite \textit{cualquiera} ‘any’ to a predicate modifier \textit{cualquiera} ‘common’ is the first relevant meaning change. The previous changes represent (only) grammatical changes rather than meaning changes as the free choice meaning of ‘which-ever one wants’ and ‘any’ are preserved in the grammatical change from a relative clause into an indefinite. The meaning change from ‘any’ to ‘common’ has further induced syntactic/grammatical changes from a predicate modifier into a predicative complement without nominal modification and later into a degree predicate in Argentinian Spanish.
As for the question of how much the meaning change involves compositionality (\cite{Frege1997, Eckardt2006, Aloni2021, Bar-asher-siegal2020, Condoravdi2015}), my answer is that the reanalysis process from ‘any’ to ‘common’ is based on the compositionality principle, because the derivation of the meaning of \textit{cualquiera} in the context of negation is compositional, that is, the negation of free choice alternatives leads to the assertion of a specific alternative. This step was necessary for speakers to make an abstraction of properties from the contrast of focus alternatives (unspecific vs. a specific contextual alternative). However, the compositionality itself is not what triggers the change. It is the characterisation of focus alternatives as properties, along with the high frequency of the negative context, that leads to the lexicalisation of the pragmatic context and the lexicalisation of the property with the meaning ‘common’ (for the role of frequency in the context of language change, see \cite{Pagel2007, Bybee2006}). 
Still, the contextual or enriched meaning is derivable compositionally from the use of \textit{cualquiera} in the negation context and focus association of postnominal \textit{cualquiera}. The derogatory meaning of \textit{cualquiera} as ‘nothing special’ or ‘bad’ is an instance of pejoration of interpreting ‘common’ in relation to a pragmatic scale according to which the property of being ‘common’ is interpreted as being non-distinguished or ordinary in the context of evaluation. The pejoration of the meaning ‘common/standard’ is a well-known process in other languages, e.g. Latin \textit{vulgaris} ‘common’ >{}> French \textit{vulgaire} ‘vulgar’/ ‘obscene’, etc. (\cite {Borkowska2007}).
The meaning change of the pronoun \textit{cualquiera} is based on the pragmatic strengthening of the meaning ‘anything’ in the context of verbs of saying (e.g. \textit{decir, mandar, hablar} ‘say’ in Argentinian Spanish), whereby the meaning of \textit{mandar cualquiera} ‘saying anything’ implies ‘saying anything possible’ or ‘saying unspecific things’ including contradictory or wrong things or saying random or unreflected things according to a judge (usually the speaker) in a conversation context (\cite {Grice1975, Chierchia2013} on pragmatic strengthening of indefinites). This pragmatic reasoning has been lexicalised under the meaning ‘non-sense’ and its semantic variants (e.g. ‘bullshit’) under transitive verbs such as \textit{decir/mandar} ‘say’ in Argentinian Spanish. The new meaning further generalised to other verbs, e.g. predicative verbs (e.g. ArgSp \textit{lo que me decís es cualquiera} ‘what you are telling me is bullshit/nonsensical’). The meaning change from anything to ‘nonsense’/‘bullshit’ is not unique to Argentinian Spanish, French \textit{n’importe quoi} shows similar meaning changes.


In sum, this book provides further support for the transparency and compositionality principle involved in meaning change and variation in the domain of modal indefinites like Spanish \textit{cualquiera} and it shows that {it is necessary to distinguish different uses of \textit{cualquiera} in order to understand the linguistic changes the word underwent, which have led to different grammatical functions and meanings. Some uses of \textit{cualquiera} such as postnominal and pronominal uses have acquired new (evaluative) meanings and have undergone the process of syntactic recategorisation into an adjective and a noun in Argentinian Spanish.}

